\subsection{CT3 --- External-type compression}\label{ct:3}\label{boundaried-pieces-and-external-types}
\ctsignature{\(P_{\ast\to3}\to \ctcert{C2,C3,C5};\ \ctint{P_{3\to7},P_{3\to12},P_{3\to8}};\ \ctrec{\ctnone}\).}

CT3 is an exact sequential machine for external-type compression.  Its
\texttt{Framework} fixes the ambient and branch-state types, boundaried pieces,
interfaces, compatible contexts, response semantics, the external-type
evaluator, type inclusion, replacement predicates, persistence, closure
claims, and downstream datum types.  A validated \texttt{Input} contains
\(\langle G,\mathfrak B,X\rangle\) and a certified interface for \(X\).
Finiteness, semantic equivalence, replacement exhaustion, and persistence are
produced by separate nodes.

The core plan has five independent fields:
\texttt{scope}, \texttt{equivalence}, \texttt{compression}, \texttt{defect},
and \texttt{table}.  Each field implements one node contract and depends only
on its immediate predecessor state.  The scope node returns either
\texttt{ScopeCandidate} or \texttt{ScopedState}.  The equivalence node has one
successor and constructs \texttt{EquivalenceState}; its
\texttt{EquivalenceCertificate} proves, for every compatible context, that the
stored external type decides the response of \(X\).  Failure to construct this
certificate is a pre-admission proof obligation, not a runtime alternative.

The compression node consumes that exact equivalence state.  A
\texttt{ReplacementWitness} contains an admissible smaller representative,
type inclusion, persistence of the branch condition under replacement, and the
resulting C2 claim.  The positive constructor returns a
\texttt{C2Certificate}; the negative constructor returns an
\texttt{UncompressibleState} carrying an exhaustive proof that no such witness
exists.  The defect node then returns exactly one of four indexed results: a C3
certificate, a CT7 exchange payload, a CT12 peeling payload, or a
\texttt{PersistentState}.  Only the last result reaches the finite-table node,
which returns a C5 certificate or a CT8 repeated-type payload.

Every payload is indexed by the exact uncompressible or persistent state that
justified it.  \texttt{Port} and \texttt{HandoffPlan} certify consumer
acceptance after \texttt{runCore}; they cannot change the selected terminal or
its evidence.  The theorems \texttt{run\_verified}, \texttt{run\_total}, and
\texttt{run\_trace\_valid} establish terminal soundness, totality relative to
the five supplied plans, and validity of the evidence-carrying path.

The soundness contracts are assigned to machine boundaries as follows.
\begin{itemize}
\tightlist
\item \textbf{S-Def} fixes the interface, response evaluator, external type, replacement vocabulary, and downstream datum types in \texttt{Framework} and \texttt{Input}.
\item \textbf{S-Equiv} is discharged by \texttt{EquivalenceCertificate} before replacement comparison.
\item \textbf{S-Dich} is discharged by the constructors of the scope, compression, defect, and table decisions.
\item \textbf{S-Pers} is carried by \texttt{PersistentState} and by every payload index.
\item \textbf{S-Comp} is carried by replacement exhaustion and finite-table classification.
\item \textbf{S-Rout} and \textbf{S-Trig} are discharged by the three handoff constructors and the separate consumer plan.
\end{itemize}

\begin{figure}[p]
\centering
\vspace*{-1.70cm}
\resizebox{\textwidth}{!}{%
\begin{tikzpicture}[x=1cm,y=1cm,every node/.style={font=\scriptsize}]
\node[bcwide] (entry) at (0,0) {{\tiny\ttfamily CT3.entry}\\\textbf{Validated input}\\\(G,\mathfrak B,X,\iota_X\)};
\node[bcdec,bclemma,text width=3.45cm,minimum width=4.20cm] (scope) at (0,-2.35) {{\tiny\ttfamily CT3.decide.scope}\\\textbf{Finite type index?}};
\node[bcscopefail,text width=3.25cm,minimum width=3.25cm] (scopeout) at (7.40,-2.35) {{\tiny\ttfamily CT3.terminal.scope}\\\textbf{Scope exit}\\infinite index};
\node[bcwide,bclemma] (equiv) at (0,-4.75) {{\tiny\ttfamily CT3.certify.equivalence}\\\textbf{S-Equiv certificate}\\type decides every compatible response};
\node[bcdec,bclemma,text width=3.45cm,minimum width=4.20cm] (compression) at (0,-7.20) {{\tiny\ttfamily CT3.decide.compression}\\\textbf{Smaller valid representative?}};
\node[bcclosed] (c2) at (7.40,-7.20) {{\tiny\ttfamily CT3.terminal.c2}\\\textbf{C2}};
\node[bcroutechoice,bctheorem,text width=3.20cm,minimum width=4.10cm] (defect) at (0,-10.10) {{\tiny\ttfamily CT3.decide.defect}\\\textbf{Classify certified defect}};
\node[bcclosed,text width=2.90cm,minimum width=2.90cm] (c3) at (-8.10,-13.30) {{\tiny\ttfamily CT3.terminal.c3}\\\textbf{C3}};
\node[bcroute,bcrouteneutral,bcoutputroute,text width=3.15cm,minimum width=3.15cm] (ct7) at (-3.40,-13.30) {{\tiny\ttfamily CT3.terminal.ct7}\\\textbf{CT 7 exit}\\exchange defect};
\node[bcroute,bcrouteneutral,bcoutputroute,text width=3.15cm,minimum width=3.15cm] (ct12) at (1.40,-13.30) {{\tiny\ttfamily CT3.terminal.ct12}\\\textbf{CT 12 exit}\\load and measure};
\node[bcroutechoice,bctheorem,text width=3.20cm,minimum width=4.10cm] (table) at (7.15,-13.30) {{\tiny\ttfamily CT3.decide.table}\\\textbf{Classify persistent table}};
\node[bcclosed,text width=2.90cm,minimum width=2.90cm] (c5) at (4.75,-16.45) {{\tiny\ttfamily CT3.terminal.c5}\\\textbf{C5}};
\node[bcroute,bcrouteneutral,bcoutputroute,text width=3.15cm,minimum width=3.15cm] (ct8) at (9.55,-16.45) {{\tiny\ttfamily CT3.terminal.ct8}\\\textbf{CT 8 exit}\\repeated types};
\draw[bcarr] (entry)--(scope);
\draw[bcscopearr] (scope)--node[bclabel,above]{exit}(scopeout);
\draw[bcarr] (scope)--node[bclabel,right]{ready}(equiv);
\draw[bcarr] (equiv)--node[bclabel,right]{certified}(compression);
\draw[bcarr] (compression)--node[bclabel,above]{replacement}(c2);
\draw[bcarr] (compression)--node[bclabel,right]{uncompressible}(defect);
\draw[bcarr] (defect)--node[bclabel,above,sloped]{routed account}(c3);
\draw[bchandoff] (defect)--node[bclabel,above,sloped]{\(P_{3\to7}\)}(ct7);
\draw[bchandoff] (defect)--node[bclabel,above,sloped]{\(P_{3\to12}\)}(ct12);
\draw[bcarr] (defect)--node[bclabel,above,sloped]{persistent}(table);
\draw[bcarr] (table)--node[bclabel,above,sloped]{finite table}(c5);
\draw[bchandoff] (table)--node[bclabel,above,sloped]{\(P_{3\to8}\)}(ct8);
\end{tikzpicture}%
}
\caption[Closure-tactic 3: exact external-type machine]{Closure-tactic 3 as an exact sequential machine.  Monospaced labels are the stable identifiers shared by Lean's \texttt{Graph.NodeId}, the JSON contract graph, and executable traces.  Certification nodes have one successor and represent proof obligations required before later decisions.  Every branching edge carries the indexed certificate, state, or payload named by its constructor.  The three handoff terminals retain their complete provenance and are accepted separately by \texttt{HandoffPlan}.}
\label{fig:ct3-external-type-compression}
\end{figure}
\FloatBarrier
